%------------------------------------------------------------------------------%
% Luis A. D'Afonseca
%
%------------------------------------------------------------------------------%
\documentclass[a4paper,12pt,fleqn]{article}

\usepackage{style_avaliacao}

\ShowAnswers

%------------------------------------------------------------------------------%
\begin{document}

\makeHeader{Integrais e Séries}{Prova 4 -- Turma 6}{12/12/2023}

% 01
%------------------------------------------------------------------------------%
\exercise{25}
  Escreva o polinômio de Taylor de terceiro grau, centrado no zero,
  da função seno hiperbólico
  \(
      \sinh(x)=\frac{e^x-e^{-x}}{2}
  \)

\begin{answer}
  O polinômio de Taylor de terceiro grau, centrado em zero, é
  \[
    P_3(x) = \sinh(0) + \sinh'(0) x + \frac{\sinh''(0)}{2}x^2 + \frac{\sinh^{(3)}(0)}{3!}x^3
  \]
  Calculando as derivadas
  \begin{align*}
    \sinh^{(1)}(x) & = \frac{1}{2}\frac{d}{dx}\left(e^x-e^{-x}\right)
                     = \frac{1}{2}\left(e^x+e^{-x}\right) \\[2mm]
    \sinh^{(2)}(x) & = \frac{1}{2}\frac{d}{dx}\left(e^x+e^{-x}\right)
                     = \frac{1}{2}\left(e^x-e^{-x}\right) \\[2mm]
    \sinh^{(3)}(x) & = \frac{1}{2}\frac{d}{dx}\left(e^x-e^{-x}\right)
                     = \frac{1}{2}\left(e^x+e^{-x}\right)
  \end{align*}
  Avaliando as derivadas em zero
  \begin{align*}
  \sinh^{(0)}(0) & = \frac{e^0-e^{-0}}{2} = \frac{1-1}{2} = 0  \\[2mm]
  \sinh^{(1)}(0) & = \frac{e^0+e^{-0}}{2} = \frac{1+1}{2} = 1  \\[2mm]
  \sinh^{(2)}(0) & = \frac{e^0-e^{-0}}{2} = \frac{1-1}{2} = 0  \\[2mm]
  \sinh^{(3)}(0) & = \frac{e^0+e^{-0}}{2} = \frac{1+1}{2} = 1
  \end{align*}
  Portanto
  \[
    P_3(x) = x + \frac{x^3}{6}
  \]
\end{answer}

% 02
%------------------------------------------------------------------------------%
\exercise{25}
Seja \( f(x) = \frac{2}{1+x}\)
\begin{enumerate}[label=\alph*)]
  \item\ [10] encontre um padrão para \(f^{(n)}(x)\)
  \item\ [15] usando o padrão do item anterior, encontre a
                   Serie de Taylor de $f(x)$, centrada em zero
\end{enumerate}

\begin{answer}
  \begin{enumerate}[label=\alph*)]
    \item Calculando as primeiras derivadas
    \begin{align*}
      f^{(0)}(x) & = \frac{2}{1+x} = 2(1+x)^{-1} \\[2mm]
      f^{(1)}(x) & = \frac{\,d\;}{\,dx\;}\left[2(1+x)^{-1}\right]
      = 2(-1)(1+x)^{-1-1} (1+x)'
      = -2 (1+x)^{-2} \\[2mm]
      f^{(2)}(x) & = \frac{d^2}{dx^2}\left[2(1+x)^{-1}\right]
      = \hpm 2 \cdot 2                 (1+x)^{-3} \\[2mm]
      f^{(3)}(x) & = \frac{d^3}{dx^3}\left[2(1+x)^{-1}\right]
      = -2 \cdot 2 \cdot 3         (1+x)^{-4} \\[2mm]
      f^{(4)}(x) & = \frac{d^4}{dx^4}\left[2(1+x)^{-1}\right]
      = \hpm 2 \cdot 2 \cdot 3 \cdot 4 (1+x)^{-5} \\[2mm]
      f^{(5)}(x) & = \frac{d^5}{dx^5}\left[2(1+x)^{-1}\right]
      = - 2 \cdot 2 \cdot 3 \cdot 4 \cdot 5 (1+x)^{-6}
    \end{align*}
    Por inspeção, vemos que
    \[
      f^{(n)}(x)
      = \frac{d^n}{dx^n}\left[2(1+x)^{-1}\right]
      = 2(-1)^n n! (1+x)^{-n-1}
      = \frac{2 (-1)^n n!}{(1+x)^{n+1}}
    \]
    \item Os coeficientes da série de Taylor são
    \[
      c_n
      = \frac{f^{(n)}(0)}{n!}
      = \frac{2(-1)^n n!}{(1-0)^{n+1}n!}
      = 2(-1)^n
    \]
    Montando a série temos
    \[
      T(x)
      = \sum_{n=0}^\infty c_n x^n
      = \sum_{n=0}^\infty 2(-1)^n x^n
    \]
  \end{enumerate}
\end{answer}


% 03
%------------------------------------------------------------------------------%
\exercise{25}
Considerando a série de potências
\(
  \sum_{n=0}^\infty (-2)^n x^n
\)
\begin{enumerate}[label=\alph*)]
  \item\ [15] determine seu raio de convergência;
  \item\ [10] calcule sua primitiva.
\end{enumerate}

\begin{answer}
\begin{enumerate}[label=\alph*)]
\item
  \textbf{Solução 1}
  Reescrevendo a série temos
  \[
    T(x)
    = \sum_{n=0}^\infty (-2)^n x^n
    = \sum_{n=0}^\infty (-2x)^n
  \]
  que é uma série geométrica com \(\alpha = 1\) e \(r = -2x\),
  portanto, convergente apenas quando
  \begin{align*}
    \abs{r} & < 1 \\
    \abs{-2x} & < 1 \\
    2\abs{x} & < 1 \\
    \abs{x} & < \frac{1}{2}
  \end{align*}
  Portando o raio de convergência é \(R=\sfrac{1}{2}\)

  \textbf{Solução 2}
  Usando o teste da razão
  \begin{align*}
    a_n     & = (-2)^n x^n         & \abs{a_n}     & = 2^n\abs{x}^n         \\
    a_{n+1} & = (-2)^{n+1} x^{n+1} & \abs{a_{n+1}} & = 2^{n+1}\abs{x}^{n+1}
  \end{align*}
  \begin{align*}
    \rho
    & = \lim_{n\to\infty} \frac{\abs{a_{n+1}}}{\abs{a_n}} \\[2mm]
    & = \lim_{n\to\infty} \frac{2^{n+1}\abs{x}^{n+1}}{2^n\abs{x}^n} \\[2mm]
    & = \lim_{n\to\infty} 2\abs{x} \\[2mm]
    & = 2\abs{x}
  \end{align*}
  A série converge quando \(\rho<1\), ou seja, \(\abs{x}<\sfrac{1}{2}\),
  portanto, o raio é  \(R=\sfrac{1}{2}\)

\item
  \begin{align*}
    F(x)
    & = \int \sum_{n=0}^\infty (-2)^n x^n \,dx \\[2mm]
    & = \sum_{n=0}^\infty (-2)^n \int x^n \,dx \\[2mm]
    & = C + \sum_{n=0}^\infty (-2)^n \frac{x^{n+1}}{n+1} \\[2mm]
    & = C + \sum_{n=0}^\infty \frac{(-2)^n}{n+1}x^{n+1}
  \end{align*}
\end{enumerate}
\end{answer}



% 04
%------------------------------------------------------------------------------%
\exercise{25}
Escreva uma aproximação para \(\cos\left(\sfrac{\pi}{8}\right)\)
utilizando o polinômio de Taylor de quarto grau e
estime o erro dessa aproximação.
Não é necessário encontrar a representação decimal dos valores numéricos.

\begin{answer}
  Polinômio de Taylor
  \[
    P_4(x) = \sum_{n=0}^{4} \frac{f^{(n)}(0)}{n!} x^n
  \]
  Calculando as derivadas do cosseno e as avaliando em zero
  \begin{align*}
    \cos^{(1)}(x) & = -\sin(x) &
    \cos^{(1)}(0) & = -\sin(0) = 0 \\
    \cos^{(2)}(x) & = -\cos(x) &
    \cos^{(2)}(0) & = -\cos(0) = -1 \\
    \cos^{(3)}(x) & =  \sin(x) &
    \cos^{(3)}(0) & =  \sin(0) = 0 \\
    \cos^{(4)}(x) & =  \cos(x) &
    \cos^{(4)}(0) & =  \cos(0) = 1
  \end{align*}
  então
  \begin{align*}
    P_4(x)
    & = \cos(0)
    - \sin(0)\, x
    - \cos(0)\, x^2
    + \sin(0)\, x^3
    + \cos(0)\, x^4 \\
    & = 1 - x^2 + x^4
  \end{align*}
  Em \(x=\sfrac{\pi}{8}\) temos
  \[
    \cos\left(\frac{\pi}{8}\right)
    \approx
    P_4\left(\frac{\pi}{8}\right)
    = 1
    - \frac{\pi^2}{8^2}
    + \frac{\pi^4}{8^4}
  \]

  Pelo \textbf{Teorema de Taylor}, sabemos que o erro cometido ao aproximarmos
  $\cos(x)$ pelo seu polinômio de Taylor centrado em zero de grau 4 é
  \[
    R_4(x)
    = \frac{\cos^{(5)}(c)}{(5)!} \, x^5
    = -\frac{\sin(c)}{120}x^5
  \]
  para algum $c$ entre zero e $x$.
  Avaliando o erro em \(x=\frac{\pi}{8}\) temos
  \[
    \abs{R_4\left(\frac{\pi}{8}\right)}
    = \abs{-\frac{\sin(c)}{120}\left(\frac{\pi}{8}\right)^5}
    = \frac{\abs{\sin(c)}}{120}\left(\frac{\pi}{8}\right)^5
  \]
  Como para $c$ entre zero e \(\sfrac{\pi}{8}\) o máximo do seno será
  em \(\sfrac{\pi}{8}\), portanto
  \[
    \abs{R_4\left(\frac{\pi}{8}\right)}
    \leq \frac{\abs{\sin\left(\dfrac{\pi}{8}\right)}}{120}\left(\frac{\pi}{8}\right)^5
  \]
\end{answer}

%------------------------------------------------------------------------------%
\end{document}
%------------------------------------------------------------------------------%
